\documentclass[10pt]{article}

\usepackage{amsmath}
%
\usepackage{amssymb}
%
\usepackage{enumerate}
%
\usepackage[width=7in,height=10in,top=0.75in,centering]{geometry}
%
\usepackage[dvipsnames]{xcolor}
%
\usepackage{tikz}
%
\usepackage{multicol}
%
\usepackage{hyperref}
%
\usepackage{marginnote}
%
%%%%%%%%%%%%%%%% HEADER %%%%%%%%%%%%%%%%%%%%%
\usepackage{fancyhdr}
\pagestyle{fancy}
%%
%\lfoot{Left footer}
%\cfoot{Center of footer}
%\rfoot{Right of footer}
%%
\lhead{Survey of Calculus I}
\chead{Week 7 Workshop}
\rhead{\thepage}
%
%%%%%%%%%%%%%%%% HEADER %%%%%%%%%%%%%%%%%%%%%
%
%%%%%%%%%%%%%%%% PGFPLOTS %%%%%%%%%%%%%%%%%%%
\usepackage{pgfplots}
\parindent0pt

\begin{document}

Let $f(x)$ be the function represented by the following graph:

\begin{center}
\begin{tikzpicture}
\begin{axis}[
width=12cm, height=10cm,
axis lines=middle,
xmin=-4, xmax=10,
ymin=-12, ymax=12,
xlabel={$x$},
ylabel={$y=f(x)$},
xtick=\empty,
ytick=\empty,
samples=100,
restrict y to domain=-20:20,
enlargelimits=false,
legend pos=north west,
legend cell align={left}
]

\addplot[black, thick, domain=-4:3.9] {-((x-2)^3)/(x-4) + 2};

\addplot[black, thick, domain=4.1:8] {((x-9)*(x-6)^2)/(-(x-4))};

\addplot[black, thick, domain=8:10] {(x-8)^3 + 1};


\draw[dashed, black, thick] (axis cs:4,-15) -- (axis cs:4,15);

\draw[dashed, black, thick] (axis cs:2,0) -- (axis cs:2,2);
\draw[dashed, black, thick] (axis cs:0,2) -- (axis cs:2,2);
\filldraw[black] (axis cs:2,2) circle (1.5pt);
\node[below] at (axis cs:2,0) {$2$};
\node[left] at (axis cs:0,2) {$2$};

\node[below] at (axis cs:4.2,0) {$4$};
\node[below] at (axis cs:6,0) {$6$};
\node[below] at (axis cs:8,0) {$8$};

\draw[dashed, black, thick] (axis cs:8,0) -- (axis cs:8,1);
\draw[dashed, black, thick] (axis cs:0,1) -- (axis cs:8,1);
        % Open circle for the limit from the left
\filldraw[black, fill=white] (axis cs:8,1) circle (1.5pt);
        % Solid dot for the defined value at x=8
        \filldraw[black] (axis cs:8,0) circle (1.5pt);
\node[left] at (axis cs:0,1) {$1$};
 \end{axis}
\end{tikzpicture}
\end{center}

\begin{enumerate}[(a)]
    \item Find all the critical numbers of $f$.
    \vfill 
    \item At which critical number(s) does $f$ attain local extrema? Find the intervals where $f$ is increasing or decreasing.
    \vfill
    \item Find all the inflection points of $f$.
    \vfill
\end{enumerate}

\textbf{Bonus problem:} Let us define a new function $g(x) = (f(x))^2 + 1$ that depends on the above function $f$.
\begin{enumerate}[(a)]
    \item Find all the critical numbers of $g$. (Hint: Compute $g'(x)$ and determine how the critical numbers of $g(x)$ depend on $f(x)$ and $f'(x)$)
    \vfill
    \item From the inflection points of $f$ that you have found, which of them are also the inflection points of $g$? (Hint: Compute $g''(x)$. What happens if you plug the inflection points of $f$ into $g''(x)$?)
    \vfill
    \item What is the concavity of $g$ at $x = 0$ (Hint: What is the sign of $g(x)$ when you plug $x = 0$ into $g''(x)$?)
    \vfill
\end{enumerate}

\end{document}